\documentclass{beamer}

\usepackage{amsmath,amsfonts,amssymb,bm}
\usepackage{quantikz}
\usepackage{graphicx}

\usetheme{Madrid}

\title{Quantum Approximate Optimization Algorithm (QAOA)}
\subtitle{From Many-Body Physics to Optimization}
\author{}
\date{}

\begin{document}

\frame{\titlepage}

%================================================
\section{Motivation and Context}
%================================================

\begin{frame}{Why Quantum Optimization?}
\begin{itemize}
\item Many real-world problems are combinatorial
\item Classical complexity: NP-hard
\item Examples:
\begin{itemize}
\item MaxCut
\item Scheduling
\item Portfolio optimization
\end{itemize}
\item Quantum computers offer exponential Hilbert space
\end{itemize}
\end{frame}

%------------------------------------------------

\begin{frame}{Physics Perspective}
\begin{itemize}
\item Optimization $\leftrightarrow$ ground state search
\item Encode cost function as Hamiltonian
\item Solve:
\[
\hat{H}_C |x\rangle = C(x)|x\rangle
\]
\item Quantum many-body system as solver
\end{itemize}
\end{frame}

%------------------------------------------------

\begin{frame}{NISQ Era Constraints}
\begin{itemize}
\item Limited qubits
\item No error correction
\item Short circuit depth required
\item Variational algorithms:
\begin{itemize}
\item Hybrid quantum-classical
\end{itemize}
\end{itemize}
\end{frame}

%================================================
\section{QUBO and MaxCut}
%================================================

\begin{frame}{QUBO Problems}
\[
C(x) = x^T Q x
\]
\begin{itemize}
\item Binary variables $x_i \in \{0,1\}$
\item NP-hard
\item Universal form for optimization problems
\end{itemize}
\end{frame}

%------------------------------------------------

\begin{frame}{Mapping to Ising Model}
\[
z_i = 2x_i - 1
\]
\[
\hat{H}_C = \sum_{i,j} J_{ij} Z_i Z_j + \sum_i h_i Z_i
\]
\begin{itemize}
\item Enables quantum implementation
\end{itemize}
\end{frame}

%------------------------------------------------

\begin{frame}{MaxCut Problem}
\begin{itemize}
\item Partition graph into two sets
\item Maximize edge cuts
\[
C(x)=\sum_{i,j} w_{ij} x_i(1-x_j)
\]
\end{itemize}
\end{frame}

%------------------------------------------------

\begin{frame}{Figure: MaxCut}
\centering
\includegraphics[width=0.6\textwidth]{example-image}
\\
\small Placeholder: Graph partition
\end{frame}

%================================================
\section{Variational Quantum Algorithms}
%================================================

\begin{frame}{Variational Principle}
\[
E(\theta) = \langle \psi(\theta)|H|\psi(\theta)\rangle
\]
\begin{itemize}
\item Minimize via classical optimizer
\end{itemize}
\end{frame}

%------------------------------------------------

\begin{frame}{Hybrid Loop}
\begin{itemize}
\item Quantum: state preparation + measurement
\item Classical: optimization
\end{itemize}
\end{frame}

%------------------------------------------------

\begin{frame}{VQE Structure}
\[
|\psi(\theta)\rangle = U(\theta)|0\rangle
\]
\begin{itemize}
\item General ansatz
\end{itemize}
\end{frame}

%================================================
\section{Adiabatic Quantum Computing}
%================================================

\begin{frame}{Adiabatic Theorem}
\begin{itemize}
\item Slow evolution preserves ground state
\end{itemize}
\[
H(t) = (1-s)H_M + sH_C
\]
\end{frame}

%------------------------------------------------

\begin{frame}{Time Evolution Operator}
\[
U = \mathcal{T}\exp\left(-i\int H(t)dt\right)
\]
\end{frame}

%------------------------------------------------

\begin{frame}{Trotterization}
\[
U \approx \prod e^{-iH_C\Delta t}e^{-iH_M\Delta t}
\]
\end{frame}

%================================================
\section{QAOA Construction}
%================================================

\begin{frame}{QAOA Idea}
\begin{itemize}
\item Discretized adiabatic evolution
\item Variational parameters replace schedule
\end{itemize}
\end{frame}

%------------------------------------------------

\begin{frame}{QAOA Ansatz}
\[
|\psi_p\rangle =
\prod_{k=1}^p e^{-i\beta_k H_M} e^{-i\gamma_k H_C} |s\rangle
\]
\end{frame}

%------------------------------------------------

\begin{frame}{Initial State}
\[
|s\rangle = |+\rangle^{\otimes n}
\]
\end{frame}

%------------------------------------------------

\begin{frame}{Cost Hamiltonian (MaxCut)}
\[
H_C = \frac{1}{2}\sum (1 - Z_i Z_j)
\]
\end{frame}

%------------------------------------------------

\begin{frame}{Mixer Hamiltonian}
\[
H_M = \sum X_i
\]
\end{frame}

%================================================
\section{Circuit Implementation}
%================================================

\begin{frame}{QAOA Circuit (p=1)}
\[
\begin{quantikz}
\lstick{|0\rangle} & H & \gate{U_C(\gamma)} & \gate{U_M(\beta)} & \meter{} \\
\lstick{|0\rangle} & H & \gate{U_C(\gamma)} & \gate{U_M(\beta)} & \meter{}
\end{quantikz}
\]
\end{frame}

%------------------------------------------------

\begin{frame}{Two-Qubit Interaction}
\[
e^{-i\gamma Z_i Z_j}
\]
\begin{quantikz}
\lstick{} & \ctrl{1} & \gate{R_Z} & \ctrl{1} & \qw \\
\lstick{} & \targ{} & \qw & \targ{} & \qw
\end{quantikz}
\end{frame}

%------------------------------------------------

\begin{frame}{Mixer Gate}
\[
e^{-i\beta X} = R_X(2\beta)
\]
\end{frame}

%================================================
\section{p=1 Analytical Solution}
%================================================

\begin{frame}{Expectation Value}
\[
F(\gamma,\beta) = \langle \psi | H_C | \psi \rangle
\]
\end{frame}

%------------------------------------------------

\begin{frame}{Edge Contribution}
\[
\langle Z_i Z_j \rangle =
\sin(4\beta)\sin(2\gamma)
\]
\end{frame}

%------------------------------------------------

\begin{frame}{Energy Expression}
\[
F = \sum \frac{1}{2}(1 - \sin(4\beta)\sin(2\gamma))
\]
\end{frame}

%------------------------------------------------

\begin{frame}{Optimal Parameters}
\[
\beta = \frac{\pi}{8}, \quad \gamma = \frac{\pi}{4}
\]
\end{frame}

%================================================
\section{Performance and Scaling}
%================================================

\begin{frame}{Approximation Ratio}
\[
\alpha = \frac{F}{C_{\max}}
\]
\end{frame}

%------------------------------------------------

\begin{frame}{Scaling with p}
\begin{itemize}
\item Larger $p$ → better results
\item $p \to \infty$ → adiabatic limit
\end{itemize}
\end{frame}

%------------------------------------------------

\begin{frame}{Figure: Performance vs p}
\includegraphics[width=0.7\textwidth]{example-image}
\end{frame}

%================================================
\section{Optimization Landscape}
%================================================

\begin{frame}{Barren Plateaus}
\[
\text{Var}(\nabla C) \sim e^{-n}
\]
\end{frame}

%------------------------------------------------

\begin{frame}{Mitigation Strategies}
\begin{itemize}
\item Problem-inspired ansatz
\item Layerwise training
\end{itemize}
\end{frame}

%================================================
\section{Advanced Variants}
%================================================

\begin{frame}{QAOA Variants}
\begin{itemize}
\item Multi-angle QAOA
\item ADAPT-QAOA
\item Counterdiabatic QAOA
\end{itemize}
\end{frame}

%------------------------------------------------

\begin{frame}{Physical Interpretation}
\begin{itemize}
\item QAOA = quantum control protocol
\item Discrete quenches
\end{itemize}
\end{frame}

%================================================
\section{Noise and Hardware}
%================================================

\begin{frame}{Noise Effects}
\begin{itemize}
\item Decoherence
\item Gate errors
\end{itemize}
\end{frame}

%------------------------------------------------

\begin{frame}{Mitigation}
\begin{itemize}
\item Error mitigation
\item Shallow circuits
\end{itemize}
\end{frame}

%================================================
\section{Outlook}
%================================================

\begin{frame}{Quantum Advantage?}
\begin{itemize}
\item Still open question
\item Depends on problem structure
\end{itemize}
\end{frame}

%------------------------------------------------

\begin{frame}{Future Directions}
\begin{itemize}
\item Better ansätze
\item Integration with ML
\item Hardware co-design
\end{itemize}
\end{frame}

%------------------------------------------------

\begin{frame}{Summary}
\begin{itemize}
\item QAOA connects:
\begin{itemize}
\item Many-body physics
\item Optimization theory
\item Quantum control
\end{itemize}
\item Promising NISQ algorithm
\end{itemize}
\end{frame}

\end{document}
